% Part of main.tex; compile main.tex.
\subsection{Compute surrogates}
\label{sec:surrogates}
\begin{table}[t]
\centering
\small
\setlength{\tabcolsep}{3pt}
\resizebox{\columnwidth}{!}{%
\begin{tabular}{@{}l l l r@{}}
\toprule
Surrogate & Resource exercised & Implementation & Lines \\
\midrule
sleep (baseline) & none & \texttt{nanosleep} & \clSleep \\
spin & one core (ALU) & xorshift loop & \clSpin \\
DGEMM & CPU, caches & OpenBLAS, $256^2$, 1 thread & \clDgemm \\
STREAM & memory bandwidth & triad, $3\times 32$\,MiB/rank & \clStream \\
halo & MPI, memory copies & 3-D 6-face exchange & \clHalo \\
heat & CPU, memory, MPI & 7-pt Jacobi, $128^3$/rank & \clHeat \\
GPU, host blocks & GPU only & FMA kernel, blocking sync & \clGpuKernel \\
GPU, host spins & GPU, one core & FMA kernel, spin sync & \clGpuKernel \\
\bottomrule
\end{tabular}}
\caption{Compute surrogates used in place of the computation phase. All are
calibrated to the same standalone duration $t_{comp}$. Lines refer to
Listing~\ref{lst:kernels} (CPU) and Listing~\ref{lst:gpu} (GPU).}
\label{tab:surrogates}
\end{table}

Table~\ref{tab:surrogates} lists the compute surrogates. Each surrogate
performs a fixed amount of \emph{work} rather than running for a fixed
\emph{time}. A busy loop bounded by time, similar to \texttt{sleep}, would end
on schedule no matter how much CPU time it received, and therefore could not
reveal interference. Before the measured steps, the work of each surrogate is
calibrated with all ranks running concurrently and no I/O in flight, so that
its standalone duration matches $t_{comp}$. The calibration is shown in
\srcCalibrate, and it is timed by \srcTimedMax. Because the work is linear in
the number of units, the calibration targets at most 2\,s and scales the
result to the full $t_{comp}$. The standalone time $t_{comp}^{alone}$ is then
measured directly (\srcBaseline). As a result, differences in run time
reflect interference rather than a change in the nominal compute length.

The STREAM triad~\cite{mccalpin1995stream} uses arrays that together exceed
the L3 cache. The heat mini-app is a 3-D Jacobi stencil with halo exchange,
and it stands in for a real simulation. Its sweep is shown in \srcHeat, and
its ghost-face exchange in \srcHeatHalo. The halo surrogate reuses the same
six-neighbour exchange (\srcExchange) on a 3-D Cartesian decomposition
(\srcCart). For the GPU surrogates, all ranks share the single GPU, and the
output field is copied to the host at the end of each compute phase
(\srcGpuRun). The wait mode of the host thread is selected in \srcGpuSetup.
Fig.~\ref{fig:loop} shows the loop that all configurations share. The
complete source code of the benchmark, its compute kernels, I/O backends, run
scripts, library patches, and analysis scripts is reproduced in
Appendix~\ref{app:code}.

\begin{figure}[t]
\begin{lstlisting}[language=C,numbers=none,xleftmargin=0pt]
calibrate units so that kernel(units) takes t_comp
  (all ranks concurrently, no I/O); measure t_alone
backend->init()
for (s = 0; s < steps; s++) {
  a = now();
  backend->compute_begin();      /* ADIOS2 Guided hint */
  kernel_run(k, units);          /* fixed work, not time */
  kernel_output(k, buf);         /* GPU: D2H copy      */
  backend->compute_end();
  b = now();
  backend->issue(s, buf);        /* sync: write + wait;
                    async: wait(s-1), then enqueue(s) */
  MPI_Barrier();                 /* I/O phase ends on all ranks */
  t_comp[s] = b - a;  t_io[s] = now() - b;
}
backend->wait(); backend->fini();  /* last step drains */
\end{lstlisting}
\caption{The epoch loop of \texttt{aio\_bench}, simplified from
\srcLoop in the appendix. At most one asynchronous step is
in flight, as in h5bench.}
\label{fig:loop}
\end{figure}

\subsection{Workload}
\label{sec:workload}
The I/O pattern follows the write kernel of h5bench and VPIC-IO. In each
timestep, every rank writes a contiguous block of a shared 1-D array of
doubles. We vary the block from 16\,MiB to 1\,GiB per rank, which is 0.75\,GiB
to 48\,GiB per timestep with 48 ranks, and we write 10 timesteps. As a result,
a single run writes between 7.5\,GiB and 480\,GiB. The largest size is
measured only for the two fastest backends, POSIX and ADIOS2, because MPI-IO
and HDF5 need 30 to 40 minutes per run at that size. The main experiments
use 256\,MiB per rank, which is 12\,GiB per step and 120\,GiB per run.

Unless stated otherwise, $t_{comp}$ is set to $1.25\times$ the synchronous
I/O time per step of an I/O-only run with the same backend, size, rank count,
and number of steps. This I/O-only run is measured at the start of each
experiment (\srcIoBaseline). With this choice, a sleep-based run can hide
the I/O completely, which is the configuration used in prior
work~\cite{tang2022async,ravi2023evaluating}. It is important to note that
the ratio achieved in the experiments differs from 1.25, because the I/O
time varies between runs. We therefore report the measured ratio for each
backend in Section~\ref{sec:rq3}. In addition, each step ends with a barrier.
Without it, ranks that finish their I/O early would charge the imbalance to
the next compute phase, because communicating kernels wait for their
neighbours.

Each configuration is run once. The main experiment is defined in
\srcPhaseMain, the size sweep in \srcPhaseSize and \srcPhaseSizeLarge, and
each configuration is executed by \srcRunOne. Repeated runs are left to
future work (Section~\ref{sec:future}). The only configurations that we ran
twice are the I/O-only baselines, and they show how large the run-to-run
variation is on this node. Across \VarPairs pairs, the two runs differ by a
median of \VarMedian\%, by more than \VarPseventyfive\% for a quarter of
them, and by up to \VarMax\%. This variation is mostly due to the different
state of the page cache and of write-back between runs.
Table~\ref{tab:codemap} in the appendix maps every mechanism described in
this section to its implementation.

\subsection{Metrics}
\label{sec:metrics}
We use the following metrics, where $n$ is the number of steps:
\begin{itemize}
  \item \textbf{Run time relative to the sleep prediction.} This metric
        directly measures how much longer the application runs than a
        sleep-based evaluation predicts. For a surrogate $k$, the sleep run
        predicts an asynchronous run time of
        $T_{async}(\texttt{sleep}) + n\,(t_{comp}^{alone}(k) -
        t_{comp}^{alone}(\texttt{sleep}))$, which corrects for small
        differences in the calibrated compute length. We divide the
        observed asynchronous run time with $k$ by this prediction
        (\srcExcess). Because
        it compares asynchronous runs only, this metric does not depend on
        the synchronous runs, whose I/O time varies between runs.
  \item \textbf{End-to-end speedup} $S = T_{sync}/T_{async}$ over all
        timesteps. This is the metric most prior work reports.
  \item \textbf{Cost attribution.} We split the asynchronous run time into a
        compute part $n\,t_{comp}^{\text{async}}$ and an exposed-I/O part
        $T_{async} - n\,t_{comp}^{\text{async}}$. The exposed-I/O part
        includes the waits for background I/O, the issue overhead, and the
        final drain. We further split the compute part into the compute time
        of the synchronous run, $n\,t_{comp}^{\text{sync}}$, and the
        inflation, $n\,(t_{comp}^{\text{async}} - t_{comp}^{\text{sync}})$.
        With perfect overlap and no interference, $T_{async}$ would equal
        $n\,t_{comp}^{\text{sync}}$ plus one step of I/O. Any excess is
        charged to the compute side (inflation) or to the I/O side (exposed
        I/O).
  \item \textbf{Compute slowdown}
        $t_{comp}^{\text{async}} / t_{comp}^{\text{sync}}$. This is the
        compute phase of the asynchronous run relative to that of the
        synchronous run with the same configuration, in which no I/O runs
        during the computation. We prefer this metric to the standalone
        baseline $t_{comp}^{\text{alone}}$, because the turbo frequency of
        the cores depends on whether they idle between compute phases, which
        shifts the standalone time by up to about 6\% in our measurements.
  \item \textbf{Exposed I/O}, which is the I/O time the application still
        observes per step in asynchronous mode, relative to the synchronous
        I/O time.
\end{itemize}
The per-step timings are reduced over ranks in \srcReduce, and the metrics
are computed from the resulting CSV records in \srcMetrics.
